\documentclass{article}
\usepackage[utf8]{inputenc}
\usepackage{ latexsym }
\usepackage{graphicx}
\graphicspath{ {../Figures/} }

\usepackage{amsmath}
\usepackage[a4paper, total={6.2in, 8in}]{geometry}

\usepackage[usenames,dvipsnames]{color}
\definecolor{darkblue}{rgb}{0,0,.6}
\definecolor{darkred}{rgb}{.7,0,0}
\definecolor{darkgreen}{rgb}{0,.6,0}
\definecolor{red}{rgb}{.98,0,0}
\usepackage[colorlinks,pagebackref,pdfusetitle,urlcolor=darkblue,citecolor=darkblue,linkcolor=darkred,bookmarksnumbered,plainpages=false]{hyperref}
\usepackage{amsthm}
\newtheorem*{definition}{Definition}

\title{CSCI 650 - Algorithms and Computability \\ Assignment 3}
\author{ }

\date{}

\begin{document}

\maketitle


\noindent Solutions to the written questions on this assignment should be submitted via PDF to Canvas before \textbf{October 4th at 11:59 pm}. Your Turing machine simulator and your implementation of the CYK algorithm should be submitted to INGInious by the same deadline. Make sure to justify your answers. \\

\noindent I highly encourage you to collaborate with one another. However, you must write up your own solutions \textbf{independently}. Feel free to communicate via \href{https://discord.gg/xfBWWkbm2P}{Discord} and to post questions on the appropriate forum in Canvas. Do not post solutions.\\

\noindent Have fun!


\section*{Problems}


\begin{enumerate}

    \item In this question we will write a simple function to simulate a Turing machine and use it to investigate $BB(2)$, the second busy beaver number.
    \begin{enumerate}
        \item (10 pts) Write a function \verb|simulateTM(delta, A, w)| which takes a transition function, a set of accepting states, and a string as input. This function should return a string associated with the final tape for the Turing machine along with a boolean indicating whether or not \verb|w| is accepted. We will assume that $A$ is the start state and that $b$ is the blank symbol. Submit your code on INGInious and include a code snippet of this function in your PDF submission. \verb|turingMachineTests.py|, a script meant to help with local testing, is available on Canvas.
        \item (15 pts) Consider the \href{https://en.wikipedia.org/wiki/Busy_beaver}{busy beaver problem} as discussed in class. Write code to consider all Turing machines with two states, a binary input alphabet, and a tape initially filled with 0s. Augment your code from part $(a)$ to count the number of steps each machine takes before halting up to a maximum of 100 steps. Create a table similar to the one below indicating the number and fraction of machines that take $i$ steps to finish. Include a snippet of the code you used to find these values.

        \begin{center}
            \hspace*{-1cm}\begin{tabular}{c||c|c|c|c|c|c|c|c}
                Steps & 0 & 1 & 2 & 3 & 4 & 5 & 6 & $\geq$ 100  \\ \hline
                Number & 3 & 8 & 4 & 12 & 42 & 112 & 57 & 239 \\ \hline
                Fraction & $0.1$ & $0.04$ & $0.2$ & $0.07$ & $0.06$ & $0.03$ & $0.05$ & $0.45$
            \end{tabular}
        \end{center}

        
        \item ($BB(n)$ pts for $n \geq 6$) Extra credit - Find a Turing machine associated with $BB(n)$ and prove that it is correct.
    \end{enumerate}

    \item (10 pts) Context-free grammars play an important role in many areas of computer science including in programming languages and compilers, natural language processing, and parsing among others.

    Write a function \verb|cyk(G, w)| that takes a context-free grammar and a string as input and returns a boolean indicating whether or not the string is in the language associated with the context-free grammar along with a table filled in according to the CYK algorithm as discussed in class. Submit your code on INGInious and include a code snippet of this function in your PDF submission. \verb|cykTests.py|, a script meant to help with local testing, is available on Canvas.

    \item Recall the definitions of the vertex cover and independent set problems.
    \begin{definition}[Vertex Cover]
    Given an undirected graph $G=(V,E)$, a \textbf{vertex cover} of $G$ is a subset of vertices $C \subseteq V$ such that, for every edge $\{u,v\} \in E$, either $u \in C$ or $v \in C$.
    \end{definition}
    \noindent \textbf{The Vertex Cover Problem} \\
    \noindent \textbf{Input:} An undirected graph $G=(V,E)$ and an integer $k$. \\
    \noindent \textbf{Output:} True if and only if $G$ has a vertex cover of size $k$ or less.
    \begin{definition}[Independent Set]
    Given an undirected graph $G=(V,E)$, an \textbf{independent set} is a subset of vertices $I \subseteq V$ such that no two elements of $I$ are adjacent. In particular, if $u,v \in I$, then $\{u,v\} \not \in E$.
    \end{definition}
    \noindent \textbf{The Independent Set Problem} \\
    \noindent \textbf{Input:} An undirected graph $G=(V,E)$ and an integer $k$.\\
    \noindent \textbf{Output:} True if and only if $G$ contains an independent set of size $k$ or more.
    \begin{enumerate}
        \item (10 pts) Write pseudocode (using Python syntax) for a function \verb|checkVertexCover(G, k, C)| that returns \verb|True| if $C$ is a vertex cover of $G$ of size $k$ or less and \verb|False| otherwise. Analyze the asymptotic runtime of this function.
        \item (10 pts) Write pseudocode (using Python syntax) for a function \verb|findIndependentSet(G)| that finds an independent set of maximum size in $G$. You may assume access to a function \verb|findVertexCover(G)| that is guaranteed to return a vertex cover of $G$ of minimum size in polynomial time. In particular, \verb|findIndependentSet(G)| must convert an arbitrary independent set instance to an instance of vertex cover. Analyze the asymptotic runtime of this function.
    \end{enumerate}

    \item The set cover problem is defined as follows.
    \begin{definition}[Set Cover]
    Given a set of elements $U$ (often called the universe) and subsets $S_1, S_2, \dots, S_m \subseteq U$, $C \subseteq \{1, 2, \dots, m\}$ is a \textbf{set cover} of $U$ if $\bigcup_{i \in C} S_i = U$.
    \end{definition}
    \noindent \textbf{The Set Cover Problem} \\
    \noindent \textbf{Input:} A universe $U$, sets $S_1, S_2, \dots, S_m \subseteq U$, and an integer $k$.\\
    \noindent \textbf{Output:} True if and only if there is a cover $C$ of $U$ such that $|C| \leq k$.
    \begin{enumerate}
        \item (10 pts) Write pseudocode (using Python syntax) for a function \verb|checkSetCover(U, S, k, C)| that returns \verb|True| if $C$ is a set cover of $U$ of size $k$ or less and \verb|False| otherwise. Analyze the asymptotic runtime of this function.
        \item (10 pts) Write pseudocode (using Python syntax) for a function \verb|findVertexCover(G)| that finds a vertex cover of minimum size in $G$. You may assume access to a function \verb|findSetCover(U, S)| that is guaranteed to return a minimum set cover of $U$ in polynomial time. In particular, \verb|findVertexCover(G)| must convert an arbitrary vertex cover instance to an instance of set cover. Analyze the asymptotic runtime of this function. \\
    \end{enumerate}
\end{enumerate}


\end{document}


